Die Tabellen enthalten bewusst sehr einfache Werte. Sie prüfen normale Zellen, vertikale und horizontale Verbindungen sowie eine kombinierte Verbindung beim Export von DOCX nach LaTeX.

\section{Test 1 Normale Tabelle 2 mal 4}\label{test-1-normale-tabelle-2-mal-4}

Zwei Zeilen und vier Spalten ohne verbundene Zellen.

{\def\LTcaptype{none} % do not increment counter
\begin{longtable}[]{@{}
  >{\centering\arraybackslash}p{(\linewidth - 6\tabcolsep) * \real{0.2500}}
  >{\raggedright\arraybackslash}p{(\linewidth - 6\tabcolsep) * \real{0.2500}}
  >{\raggedright\arraybackslash}p{(\linewidth - 6\tabcolsep) * \real{0.2500}}
  >{\raggedright\arraybackslash}p{(\linewidth - 6\tabcolsep) * \real{0.2500}}@{}}
\toprule\noalign{}
\begin{minipage}[b]{\linewidth}\centering
A1
\end{minipage} & \begin{minipage}[b]{\linewidth}\centering
B1
\end{minipage} & \begin{minipage}[b]{\linewidth}\centering
C1
\end{minipage} & \begin{minipage}[b]{\linewidth}\centering
D1
\end{minipage} \\
\midrule\noalign{}
\endhead
\bottomrule\noalign{}
\endlastfoot
A2 & B2 & C2 & D2 \\
\end{longtable}
}

\section{Test 2 Multirow Tabelle 5 mal 4}\label{test-2-multirow-tabelle-5-mal-4}

Fünf Zeilen und vier Spalten. A umfasst drei Zeilen und X umfasst zwei Zeilen.

{\def\LTcaptype{none} % do not increment counter
\begin{longtable}[]{@{}
  >{\centering\arraybackslash}p{(\linewidth - 6\tabcolsep) * \real{0.2500}}
  >{\raggedright\arraybackslash}p{(\linewidth - 6\tabcolsep) * \real{0.2500}}
  >{\raggedright\arraybackslash}p{(\linewidth - 6\tabcolsep) * \real{0.2500}}
  >{\raggedright\arraybackslash}p{(\linewidth - 6\tabcolsep) * \real{0.2500}}@{}}
\toprule\noalign{}
\begin{minipage}[b]{\linewidth}\centering
\textbf{Spalte 1}
\end{minipage} & \begin{minipage}[b]{\linewidth}\centering
\textbf{Spalte 2}
\end{minipage} & \begin{minipage}[b]{\linewidth}\centering
\textbf{Spalte 3}
\end{minipage} & \begin{minipage}[b]{\linewidth}\centering
\textbf{Spalte 4}
\end{minipage} \\
\midrule\noalign{}
\endhead
\bottomrule\noalign{}
\endlastfoot
\multirow{3}{=}{\centering\arraybackslash \begin{minipage}[t]{\linewidth}\centering
\textbf{A\\
3 Zeilen}\strut
\end{minipage}} & B1 & C1 & D1 \\
& B2 & C2 & D2 \\
& B3 & \multirow{2}{=}{\begin{minipage}[t]{\linewidth}\centering
\textbf{X\\
2 Zeilen}\strut
\end{minipage}} & D3 \\
B & B4 & & D4 \\
\end{longtable}
}

\section{Test 3 Multicolumn Tabelle}\label{test-3-multicolumn-tabelle}

Horizontale Verbindungen als Gegenprobe für LaTeX multicolumn.

{\def\LTcaptype{none} % do not increment counter
\begin{longtable}[]{@{}
  >{\centering\arraybackslash}p{(\linewidth - 6\tabcolsep) * \real{0.2500}}
  >{\centering\arraybackslash}p{(\linewidth - 6\tabcolsep) * \real{0.2500}}
  >{\raggedright\arraybackslash}p{(\linewidth - 6\tabcolsep) * \real{0.2500}}
  >{\raggedright\arraybackslash}p{(\linewidth - 6\tabcolsep) * \real{0.2500}}@{}}
\toprule\noalign{}
\multicolumn{2}{@{}>{\centering\arraybackslash}p{(\linewidth - 6\tabcolsep) * \real{0.5000} + 2\tabcolsep}}{%
\begin{minipage}[b]{\linewidth}\centering
\textbf{AB 2 Spalten}
\end{minipage}} & \multicolumn{2}{>{\centering\arraybackslash}p{(\linewidth - 6\tabcolsep) * \real{0.5000} + 2\tabcolsep}@{}}{%
\begin{minipage}[b]{\linewidth}\centering
\textbf{CD 2 Spalten}
\end{minipage}} \\
\midrule\noalign{}
\endhead
\bottomrule\noalign{}
\endlastfoot
A2 & B2 & C2 & D2 \\
A3 & B3 & C3 & D3 \\
A4 & B4 & C4 & D4 \\
\end{longtable}
}

\section{Test 4 Kombinierter Block}\label{test-4-kombinierter-block}

Eine Zelle umfasst gleichzeitig zwei Zeilen und zwei Spalten.

{\def\LTcaptype{none} % do not increment counter
\begin{longtable}[]{@{}
  >{\centering\arraybackslash}p{(\linewidth - 6\tabcolsep) * \real{0.2500}}
  >{\raggedright\arraybackslash}p{(\linewidth - 6\tabcolsep) * \real{0.2500}}
  >{\raggedright\arraybackslash}p{(\linewidth - 6\tabcolsep) * \real{0.2500}}
  >{\raggedright\arraybackslash}p{(\linewidth - 6\tabcolsep) * \real{0.2500}}@{}}
\toprule\noalign{}
\begin{minipage}[b]{\linewidth}\centering
\textbf{Spalte 1}
\end{minipage} & \begin{minipage}[b]{\linewidth}\centering
\textbf{Spalte 2}
\end{minipage} & \begin{minipage}[b]{\linewidth}\centering
\textbf{Spalte 3}
\end{minipage} & \begin{minipage}[b]{\linewidth}\centering
\textbf{Spalte 4}
\end{minipage} \\
\midrule\noalign{}
\endhead
\bottomrule\noalign{}
\endlastfoot
\multicolumn{2}{@{}>{\centering\arraybackslash}p{(\linewidth - 6\tabcolsep) * \real{0.5000} + 2\tabcolsep}}{%
\multirow{2}{=}{\centering\arraybackslash \begin{minipage}[t]{\linewidth}\centering
\textbf{Block\\
2 mal 2}\strut
\end{minipage}}} & C2 & D2 \\
& & C3 & D3 \\
A4 & B4 & C4 & D4 \\
\end{longtable}
}

\section{Test 5 Multirow am Tabellenrand}\label{test-5-multirow-am-tabellenrand}

Die letzte Spalte umfasst die letzten drei Zeilen und enthält einen manuellen Zeilenumbruch.

{\def\LTcaptype{none} % do not increment counter
\begin{longtable}[]{@{}
  >{\centering\arraybackslash}p{(\linewidth - 4\tabcolsep) * \real{0.3333}}
  >{\raggedright\arraybackslash}p{(\linewidth - 4\tabcolsep) * \real{0.3333}}
  >{\raggedright\arraybackslash}p{(\linewidth - 4\tabcolsep) * \real{0.3333}}@{}}
\toprule\noalign{}
\begin{minipage}[b]{\linewidth}\centering
\textbf{Spalte 1}
\end{minipage} & \begin{minipage}[b]{\linewidth}\centering
\textbf{Spalte 2}
\end{minipage} & \begin{minipage}[b]{\linewidth}\centering
\textbf{Spalte 3}
\end{minipage} \\
\midrule\noalign{}
\endhead
\bottomrule\noalign{}
\endlastfoot
A1 & B1 & \multirow{3}{=}{\begin{minipage}[t]{\linewidth}\centering
\textbf{Rand\\
3 Zeilen}\strut
\end{minipage}} \\
A2 & B2 \\
A3 & B3 \\
\end{longtable}
}
